\scanpage{28}
\mspage{24}
% archive sheet number: 27

Donc, si à chaque strate \inserted{$X$} on associe son support interne $|X|^{0}\subset I$, avec la relation d'ordre induite par $I$ (NB $|X|^{0}$ peut être vide \inserted{si $X\in\unsure{M_1|\mathcal M_1}$} !), alors $J$ est la somme \deleted{ordonnée des} \inserted{lexicographique} des $|X|^{0}$ $(X\in\unsure{\Phi|\varphi})$, $X$ suivant \unsure{dans} l'ordre $\leq$ de $M$. D'autre part, \deleted{$\Omega_0$ est vide ou un \unsure{point}}
% slanted
\margin{cette description tjrs valable}
\[
\Omega_0=
\left\{\begin{aligned}
&\varnothing\text{ si le premier élément de }\Phi\text{ n'a pas d'origine}\\
&\quad\text{i.e. est un }I_{<a},\text{ ou s'il a}\\
&\deleted{\operatorname{or}(\Phi)\cap\Phi}\quad\text{une origine }a\deleted{\unsure{\Phi}}\text{ et }a\notin\Phi\\
&a=\operatorname{or}(\Phi)\text{ sinon}
\end{aligned}\right.
\]
\[
\deleted{\Omega_1=\operatorname{ex}(\Phi)\cap}
\]
$\Omega_1$ défini dualement

\uline{NB} Il faut définir l'origine d'une préfigure, quand elle existe (i.e. quand sa première strate \deleted{\unsure{est un}} n'est pas un $I_{<a}$), et dualement, l'extrémité.

\begin{center}\rule{2.8cm}{0.4pt}\end{center}

Cas $M$ \uline{divisible}. Dans le cas $(I,\omega_0,\omega_1)$ cela signifie que pour $a<b$ $(a,b\in\unsure{J|I})$ $\exists c$ avec $a<c<b$. Dans le cas des $(I,\omega_0)$ resp. $(I,\omega_1)$, il faut ajouter que $\forall a\in I$, existe $b\in I$ t.q. $b<a$ resp $b>a$. Enfin, dans le cas de l'intervalle $I$ biouvert, il faut \inserted{alors} ajouter encore : $I\neq\varnothing$. (\uline{NB} pour dire que l'\uline{ensemble ordonné} $I$ est divisible : donc $\varnothing$ ne sera pas considéré comme ens. ordonné divisible, bien que $M(\varnothing)=\{\varnothing\}$, magasin ponctuel, le soit !)

Si $I$ est divisible, alors l'ens $M_c\subset M$ des strates ``compactes'' (i.e. de la forme $a$ et $I_{ab}$) est très dense.
